\section{Discussion}\label{sec:discussion}

We extended the refined large $N$ duality of KN to framed knots and tested the positivity of the spin form and the genus form of the refined LMOV invariants.
With defined inputs (torus knots and algebraic iterated torus knots) no negative term was found in any case examined,
up to three boxes for cables, up to nine boxes for torus knots, and over ranges of framings.
With homological or conjectural inputs, positivity is a theorem for $\yng(1)$ and is proved for families of twist and double twist knots for two boxes.
We close with open problems.

\paragraph{Positivity for all framings.}
For individual knots the exponential-polynomial structure in $x$ and the Newton criterion reduce positivity for all $x$ to finitely many checks (Theorem~\ref{thm:two-all-x}).
This is a knot-by-knot verification and not a theorem for general knots.
It is open whether the positivity of the Newton coefficients can be derived from properties of the input such as non-negativity and $\sigma$-invariance of homology.
Certificates are also missing for twist knots with $p$ and $x$ of opposite signs and for three boxes at all $x$.
We did not compare the framed unknot with the refined open BPS invariants of toric branes of Cheng and Su\l{}kowski \cite{Cheng21}.

\paragraph{Physical derivation of the framing.}
We showed that the Mari\~no--Vafa sign is necessary for the Laurent property (Proposition~\ref{prop:sign}),
but we did not derive the refined shift of the flat coordinate, $V\mapsto-(q/t)^{1/2}V$, from geometry.
That $U(1)_R$ acts trivially on the string side independently of the framing also remains an assumption.

\paragraph{Non-rectangular colours for hyperbolic knots.}
Three boxes for hyperbolic knots require a definition of the refined invariant for $\yng(2,1)$. We see three candidates:
(i) a construction through refined Racah matrices (a refinement of the arborescent calculus);
(ii) an extension of the DAHA or elliptic Hall algebra definition to hyperbolic knots (the genus-two DAHA is so far limited to $\mathfrak{sl}_2$ \cite{Arthamonov15});
(iii) a computation of HOMFLY-PT homology coloured by $\yng(2,1)$ together with a rule converting it to $\brCS$.
All are beyond the scope of this paper.
For $\mathbf{4_1}$ with four boxes the colour $\yng(3,1)$ poses the same problem (\S\ref{sec:fig8-four}).
The KN candidate for $\yng(2,1)$ fits the $Z$-factor expansion of Anokhina et al.\ (Computational Result~\ref{cr:fig8-21}), but no expansion of the same form was found for $\yng(3,1)$.
The condition \eqref{eq:HF} of Heegaard--Floer type holds for the elliptic Hall algebra values and can serve as a new constraint on $\yng(3,1)$.

\paragraph{Normalisation.}
The closed formula for the monomial normalisation (Computational Result~\ref{cr:M}) was derived under the refined satellite formula (Hypothesis~\ref{hyp:sat}).
Since the DAHA and the elliptic Hall algebra do not determine the monomial, this hypothesis remains a physical assumption.
For torus knots, the form of the corollary of Mellit in the standard normalisation was verified by computation (Computational Result~\ref{cr:mellit});
an explicit identification of the operators $d_\pm$ with $Q_{m,n}$ is left open. It does not affect the positivity results.

\paragraph{Number of boxes.}
The tests for torus knots beyond three boxes are finite computations (Table~\ref{tab:higher}).
The coefficients of the genus form grow rapidly with the number of boxes, yet no negative term appears.
A proof of positivity for all numbers of boxes from the structure of the elliptic Hall algebra is open.

\paragraph{Further questions.}
The remaining opposite-sign region for double twist knots (\S\ref{sec:two}) and a structural reason for the unimodality of the two-box invariants are also open.
